\section{Results}

\subsection{Normalization and stability}

Phase~1 returned 42 complete, 80 partial, 16 out-of-scope, and one unresolved
primary mapping.  Of the 80 partial cases routed to Phase~2, two became
complete, 77 remained partial, and one became unresolved.  The final outcome
was therefore 44/139 complete mappings (31.7\%), 77 partial, 16 out-of-scope,
two unresolved, and no schema failures.  Conservative filtering admitted 48
source-cell edges into 24 unique canonical cells.  Thirteen cells had one
supporting descriptor, six had two, four had four, and one had seven.

\begin{table*}[t]
  \centering
  \small
  \begin{minipage}[t]{.37\textwidth}
    \centering
    \begin{tabular}{@{}lrr@{}}
      \toprule
      Mapping outcome & Count & Share \\
      \midrule
      Complete & 44 & 31.7\% \\
      Partial & 77 & 55.4\% \\
      Out of scope & 16 & 11.5\% \\
      Unresolved & 2 & 1.4\% \\
      \midrule
      Source--cell edges & 48 & -- \\
      Unique cells & 24 & -- \\
      \bottomrule
    \end{tabular}
  \end{minipage}\hfill
  \begin{minipage}[t]{.60\textwidth}
    \centering
    \begin{tabular}{@{}lrrrr@{}}
      \toprule
      Item cohort & Attempts & Payloads & Accepted & Cells \\
      \midrule
      Curated $N=1$ prefix & 24 & 24 & 16 & 16 \\
      Curated $N=2$ prefix & 48 & 48 & 33 & 19 \\
      Curated $N=3$ prefix & 72 & 71 & 51 & 22 \\
      Unchanged-prompt rescue & 4 & 4 & 1 & 1 \\
      Explicit-cue intervention & 2 & 2 & 2 & 1 \\
      \midrule
      Final curated & 78 & 77 & 54 & 24 \\
      Selected fixed bank & -- & -- & 44 & 24 \\
      \bottomrule
    \end{tabular}
  \end{minipage}
  \caption{Source normalization (left) and item construction (right).
  Prefix rows reuse the same first $N$ default draws after the frozen six-item
  correction/rejudgment.  Initial pre-correction coverage was 18/21/22 cells.
  Rescue/intervention cohorts have separate provenance; accepted denotes model,
  not human, validation.}
  \label{tab:normalization-results}
\end{table*}


All eight repeated annotation units matched their primary unit exactly in
status, cells, and note.  The five repeated item diagnostics likewise agreed on
all decision flags, although their free-text notes differed.  These small
repeat sets show deterministic-looking behavior for the sampled calls, not a
general reliability estimate.

\subsection{Policy and item structure}

\begin{table*}[t]
  \centering
  \small
  \begin{minipage}[t]{.49\textwidth}
    \centering
    \begin{tabular}{@{}lrrrr@{}}
      \toprule
      Candidate family & Raw & Support & Duplicate & Eligible \\
      \midrule
      Feature value & 9 & 9 & 0 & 9 \\
      Operation & 10 & 7 & 6 & 3 \\
      Interaction & 18 & 8 & 2 & 7 \\
      Development cell & 18 & 18 & 9 & 9 \\
      \midrule
      Total / classes & 55 & 42 & 17 & 28 \\
      \bottomrule
    \end{tabular}
  \end{minipage}\hfill
  \begin{minipage}[t]{.49\textwidth}
    \centering
    \begin{tabular}{@{}lrrrr@{}}
      \toprule
      Policy & KCs & Inter. & $Q$ density & KCs/item \\
      \midrule
      Factorized & 9 & 0 & .232 & 2.09 \\
      Automated & 10 & 1 & .218 & 2.18 \\
      All supported & 16 & 7 & .180 & 2.89 \\
      Oracle all-cell & 24 & 0 & .042 & 1.00 \\
      \bottomrule
    \end{tabular}
  \end{minipage}
  \caption{Development-derived candidate space (left) and frozen policy
  granularity on the 44-item bank (right).  Duplicate means activation-equivalent
  on development items; ``oracle'' includes evaluation-cell KCs unavailable to
  the selector.}
  \label{tab:policy-results}
\end{table*}


Table~\ref{tab:policy-results} shows that policy choice materially changes the
benchmark.  Factorization shares nine KCs across the 24 cells and yields a mean
row width of 2.20.  Adding interactions increases this to 14 KCs and 2.91; the
DO-question interaction is identical to the polar-question column on this
sample, so it adds no observed distinction.  Whole-cell encoding produces 24
orthogonal columns and unit-width rows.

Deterministic checks rejected no candidate.  The automated diagnostic accepted
all 45 factorized and all 120 whole-cell items.  It rejected one of 70
interaction items for suspected answer ambiguity: the negative past
perfect-progressive target placed \emph{not} after the first auxiliary, while
the diagnostic identified a marked alternative placement consistent with the
literal prompt.  This case illustrates why deterministic realization alone
does not establish a unique pedagogically acceptable answer.

\subsection{Synthetic KT checks}

\begin{table*}[t]
  \centering
  \small
  \begin{tabular}{@{}lrrrl@{}}
    \toprule
    Supplied representation & KCs &
      Probe log loss / $\Delta$ vs.\ \kstar{} [95\% CI] &
      State RMSE & Interpretation \\
    \midrule
    All merged & 1 & $+.010225\;[.009380,.011029]$ & -- & maximal coarsening \\
    Linguistic families & 6 & $+.008132\;[.007450,.008779]$ & .146300 & union merge \\
    Generator \kstar{} & 18 & $.670627$ (reference) & \textbf{.123738} & planted truth \\
    Structural split-2 & 35 & $+.003165\;[.002744,.003581]$ & .132752 & context refinement \\
    Structural split-4 & 66 & $+.005868\;[.005281,.006473]$ & .140428 & finer refinement \\
    Exact cell & 75 & $+.015039\;[.014108,.015990]$ & .163828 & no reuse \\
    \bottomrule
  \end{tabular}
  \caption{RQ2 response prediction and oracle-only prerequisite-state recovery.
  Log-loss intervals resample 1,000 learners 2,000 times.  State RMSE compares
  the inverse-linked public prediction with the oracle minimum active-KC
  mastery on 113,000 probes; all-merged was not in the frozen mastery subset.}
  \label{tab:kt-results}
\end{table*}


Every run passed Q-matrix integrity and observable/oracle separation checks.
The policy-specific simulations contained 16,200, 24,840, and 43,200 events per
seed for factorized, interaction, and whole-cell policies.  Within each dataset,
logistic regression achieved the highest mean test AUC and lowest log loss
(Table~\ref{tab:kt-results}); fixed-parameter BKT had the highest log loss.
These results establish that the generated artifacts support chronological KT
evaluation and respond consistently across seeds.  They do not rank the KC
policies: the item counts, row widths, response-generating states, and outcome
rates all change with the policy.
